\documentclass[journal]{IEEEtran}
\usepackage{cite}
\usepackage{amsmath,amssymb,amsfonts}
\usepackage{graphicx}
\usepackage{textcomp}
\usepackage{xcolor}
\usepackage[hidelinks]{hyperref}
\usepackage{algorithm}
\usepackage{algpseudocode}
\usepackage{float}
\usepackage{placeins}
\setlength{\parskip}{0pt}




\makeatletter
\newcommand{\LEOFigureSetup}{%
  \setlength{\parskip}{0pt}%
  \setlength{\abovecaptionskip}{4pt}%
  \setlength{\belowcaptionskip}{0pt}%
  \def\@IEEEfigurecaptionsepspace{\vskip4pt\relax}%
}
\newcommand{\LEOTableSetup}{%
  \setlength{\parskip}{0pt}%
  \setlength{\belowcaptionskip}{5pt}%
  \def\@IEEEtablecaptionsepspace{\vskip5pt\relax}%
}
\makeatother



% Section V: page-top floats; one full-width float per page.
\makeatletter
\newcommand{\LEOPageTopLayout}{%
  \setcounter{topnumber}{1}%
  \setcounter{dbltopnumber}{1}%
  \renewcommand{\topfraction}{0.92}%
  \renewcommand{\dbltopfraction}{0.80}%
  \renewcommand{\textfraction}{0.08}%
  \setlength{\textfloatsep}{10pt plus 2pt minus 2pt}%
  \setlength{\dbltextfloatsep}{10pt plus 2pt minus 2pt}%
  \setlength{\@fptop}{0pt}%
  \setlength{\@dblfptop}{0pt}%
}
\makeatother

\title{Joint Handover and Partial Task Offloading in LEO Satellite Edge Networks via Heterogeneous Graph-Based Parameterized Reinforcement Learning}
\author{Author Name}
\begin{document}
\maketitle

\begin{abstract}
Integrating low Earth orbit (LEO) satellites with edge computing extends computation services to areas with limited terrestrial coverage and to terminals with constrained local resources. However, satellite motion continually changes link quality, satellite visibility, and remaining service time, while multiple users compete for limited onboard communication and computing resources. For divisible tasks, a user must select a serving satellite and determine a continuous offloading ratio matched to that satellite's link, resource, and contention conditions. We therefore formulate joint serving-satellite selection and partial task offloading as a sequential decision problem with a parameterized hybrid action space. To solve this problem, we propose HAN--PDQN, which combines a heterogeneous graph attention network (HAN) with a parameterized deep Q-network (PDQN). A heterogeneous graph distinguishes user and satellite nodes as well as serving and visibility relations. Relation-specific meta-path aggregation and hierarchical attention capture established resource sharing and potential contention among users with overlapping satellite alternatives. Conditioned on this relational representation, local task information, and candidate-satellite features, PDQN generates a continuous offloading ratio for each feasible satellite action and estimates the long-term value of the resulting satellite--ratio pair. The same contention context therefore affects both how much workload is assigned to each candidate and whether that candidate is selected. HAN--PDQN consequently supports coordinated serving-satellite selection and partial task offloading in dynamic multi-user LEO satellite edge networks.Simulations across five training seeds and network sizes of 20--40 users show that, compared with PDQN, HAN--PDQN reduces average successful-task delay, energy per successful task, and the deadline violation rate by 2.16\%, 5.72\%, and 21.26\%, respectively, while increasing the task success rate by 2.12\%.
\end{abstract}
% Insert one verified quantitative result after completing the evaluation.

\begin{IEEEkeywords}
Heterogeneous graph attention, joint handover and offloading, LEO satellite edge computing, parameterized deep reinforcement learning, partial task offloading.
\end{IEEEkeywords}

\section{Introduction}
\label{sec:introduction}

Low Earth orbit (LEO) satellite networks extend communication
coverage beyond terrestrial infrastructure and provide a
platform for delivering computing services to remote users.
By equipping satellites with onboard edge servers,
computation-intensive tasks can be processed closer to their
sources, complementing terrestrial edge computing in remote
regions, maritime environments, and other areas with limited
ground infrastructure~\cite{wang2023satellitecomputing}.
These capabilities make LEO satellite edge networks
particularly relevant to terminals whose local computing
resources cannot consistently meet application requirements.

Providing sustained computing services in such networks
requires coordinated satellite association and task offloading.
Satellite motion continuously changes link quality, the set
of visible satellites, and the remaining service time,
necessitating handovers to maintain connectivity
\cite{wang2023seamless}.
Meanwhile, users attached to the same satellite share
its communication and computing resources.
For a divisible task, each user must determine both the
satellite through which it obtains remote service and the
fraction of its workload to offload, with the remaining
workload processed locally in parallel.

These decisions are closely coupled.
The selected satellite determines the transmission conditions,
available computing resources, and queueing delay experienced
by the offloaded workload.
Conversely, the offloading fraction changes resource
occupancy and therefore affects the service available to
other users.
The coupling also extends across time: unfinished workloads
persist between decision slots, while association changes
can interrupt service and require workload reassignment.
Consequently, a favorable link at the current slot does not
necessarily identify the most suitable computing destination.
Effective control must account for the joint consequences
of satellite selection and workload allocation under
time-varying connectivity and multi-user contention.

\subsection{Related Work and Research Gap}

Graph-based learning provides a useful foundation for
representing interactions in communication networks.
Graph neural networks have been investigated for scalable
radio resource management~\cite{shen2021graph}, while
heterogeneous graph representations have been applied to
handover decisions in satellite--terrestrial integrated
networks~\cite{fu2026heterographhandover}.
These approaches organize network measurements through
relations among communicating entities, enabling decisions
to incorporate information beyond an individual link.
For joint satellite association and partial offloading,
such relational information must describe both the users
currently sharing resources and those whose future decisions
may create additional contention.

A complementary line of research considers computation
offloading in satellite edge networks.
Existing studies have jointly addressed access selection
and computation offloading~\cite{yang2025access}, and have
further incorporated communication and computing resource
allocation~\cite{yang2025accessresource}.
Multi-agent reinforcement learning has also been employed
to coordinate offloading and resource allocation
\cite{li2024marl_offloading,hassan2024comappo},
including hierarchical decision mechanisms for LEO
satellite networks~\cite{lai2025hierarchical}.
Together, these studies establish the importance of
coordinating access and computing decisions in dynamic
multi-user environments.

The combination of discrete and continuous decisions
introduces a further modeling requirement.
A discrete satellite choice specifies the service
destination, whereas a continuous offloading fraction
specifies the workload assigned to that destination.
The parameterized deep Q-network (PDQN) formulation
supports this action structure by associating continuous
parameters with discrete actions~\cite{xiong2018pdqn}.
PDQN has also been investigated for computation offloading
and resource allocation in LEO satellite edge networks
\cite{yang2023leopdqn}.
In the setting considered here, each candidate satellite
has distinct transmission, computing, and contention
conditions, making its associated offloading fraction
an integral part of the decision.

The design issue addressed in this work is how to connect
heterogeneous resource-sharing relationships to the
generation and evaluation of these satellite-specific
actions.
Users attached to the same satellite create established
contention, whereas users with overlapping satellite
alternatives may create potential contention through
subsequent association decisions.
An aggregate resource indicator summarizes the current
load but does not explicitly preserve these relationships
or the user demands underlying them.
Moreover, the value of selecting a satellite depends on
the workload fraction assigned to it.
These observations motivate a relational representation
that informs both candidate-specific workload allocation
and the long-term evaluation of the resulting
satellite--fraction pairs.

\subsection{Main Contributions}

To address this problem, we develop HAN+PDQN, which
integrates a heterogeneous graph attention network (HAN)
with parameterized reinforcement learning.
HAN organizes user demand, satellite resources, and
their interactions into a relational representation.
Conditioned on this representation and the user's
task and candidate-satellite observations, PDQN generates
an offloading fraction for each satellite action and
evaluates the corresponding action--parameter pair.
The main contributions are summarized as follows:

\begin{itemize}
    \item We formulate joint satellite association and
    partial task offloading as a sequential decision
    problem with a parameterized hybrid action space.
    The formulation captures time-varying visibility,
    stochastic task arrivals, shared communication and
    computing resources, and the service effects of
    association changes and unfinished workloads.

    \item We develop a heterogeneous graph representation
    of multi-user contention in LEO satellite edge
    networks.
    Serving relations characterize established resource
    sharing, while visibility relations capture potential
    competition through common satellite alternatives.
    Hierarchical attention integrates these relations
    with neighboring-user and satellite-topology
    information to construct user representations
    for joint control.

    \item We design a shared PDQN controller that uses
    the relational representation to generate
    satellite-specific offloading fractions and
    estimate the long-term values of the corresponding
    satellite--fraction pairs.
    Feasibility-aware action selection connects the
    destination and amount of offloaded work within
    a common decision process.
\end{itemize}

The remainder of this paper is organized as follows.
Section~\ref{sec:system_model} presents the network,
communication, computation, and handover models.
Section~\ref{sec:problem_formulation} formulates the
joint optimization problem.
Section~\ref{sec:proposed_framework} develops the
HAN+PDQN framework and its training and execution
procedures.
Section~\ref{sec:performance_evaluation} presents
the simulation settings and performance evaluation.
Finally, Section~\ref{sec:conclusion} concludes the paper.


\section{System Model}
\label{sec:system_model}

\subsection{Network Model}
\label{subsec:network_model}

We consider the multi-user LEO satellite edge computing
network shown in Fig.~\ref{fig:system_architecture}, where
ground terminals can access onboard MEC servers through
their serving satellites. Specifically, each terminal has
a local processor, while each satellite provides computing
resources shared by its associated users. We denote the
sets of users and satellites by
$\mathcal{U}=\{1,\ldots,N\}$ and
$\mathcal{S}=\{1,\ldots,M\}$, respectively.
For divisible tasks, a user can offload part of its workload
to the serving MEC and process the remaining part locally.

\begin{figure}[t]
    \centering
    \includegraphics[
        trim=3.5mm 5mm 5mm 4mm,
        clip,
        width=\columnwidth
    ]{system_architecture.pdf}
    \caption{System architecture of the multi-user LEO satellite
    edge computing network. }
    \label{fig:system_architecture}
\end{figure}

To describe the network evolution, we divide time into
slots indexed by $t\in\mathcal{T}$, each lasting $\Delta t$.
As the satellites move, the distance $d_{i,s}(t)$ and
elevation angle $\theta_{i,s}(t)$ between user $i$ and
satellite $s$ vary over time. A satellite becomes available
for service when its elevation meets the minimum requirement
$\theta_{\min}$. Accordingly, user $i$ can observe the
satellite set
\begin{equation}
\mathcal{V}_i(t)=
\left\{s\in\mathcal{S}:
\theta_{i,s}(t)\geq\theta_{\min}\right\}.
\label{eq:visible_set}
\end{equation}

For each visible satellite, we use $v_{i,s}(t)$ to denote
its remaining visibility time. Although several satellites
may be visible simultaneously, a connected user maintains
one serving association. To distinguish the association
before and after a decision, we write $b_i^{-}(t)$ and
$b_i^{+}(t)$, respectively, and use $0$ for a disconnected
state. Task admission and communication resource allocation
then follow the resulting association $b_i^{+}(t)$.

In addition to user--satellite connections, the constellation
contains ISLs that provide the satellite-neighborhood
topology for HAN. We represent task migration through
logical reassignment of residual workloads between MEC
queues, together with a fixed service-interruption overhead.

\subsection{Communication Model}
\label{subsec:communication_model}

To separate simultaneous uplink requests, we adopt OFDMA
and distribute each satellite's bandwidth among its newly
admitted uploads. In slot $t$, let $\mathcal{A}_s(t)$
collect the users admitted by satellite $s$, with
$n_s(t)=|\mathcal{A}_s(t)|$. Given total uplink bandwidth
$B_s$, each user in this admission batch receives
\begin{equation}
B_{i,s}(t)=\frac{B_s}{n_s(t)},
\qquad i\in\mathcal{A}_s(t).
\label{eq:bandwidth_sharing}
\end{equation}
The allocation is renewed for each decision slot.

For the user--satellite link, we account for free-space
propagation and additional attenuation through the
large-scale power gain
\begin{equation}
g_{i,s}(t)=
\frac{G_iG_s}
{\ell_{i,s}(t)
 \left(4\pi f_c d_{i,s}(t)/c\right)^2},
\label{eq:channel_gain}
\end{equation}
where $G_i$ and $G_s$ denote the user and satellite
antenna gains, respectively, $f_c$ is the carrier frequency,
and $c$ is the speed of light. The factor $\ell_{i,s}(t)$
combines atmospheric, pointing, polarization, and
implementation losses. Under this gain model, transmitting
with radiated power $P_i$ over bandwidth $B_{i,s}(t)$
gives the uplink rate
\begin{equation}
R^{\mathrm{up}}_{i,s}(t)=
B_{i,s}(t)
\log_2\left(
1+\frac{P_i g_{i,s}(t)}
{N_0B_{i,s}(t)}
\right),
\label{eq:uplink_rate}
\end{equation}
with $N_0$ denoting the effective noise power spectral
density.

For result return, we evaluate the communication cost
using reference downlink bandwidth $B_s^{\mathrm{dl}}$
and satellite radiated power $P_s$. Assuming the same
large-scale gain for both link directions, the corresponding
reference rate can be expressed as
\begin{equation}
R^{\mathrm{dl}}_{s,i}(t)=
B_s^{\mathrm{dl}}
\log_2\left(
1+\frac{P_s g_{i,s}(t)}
{N_0B_s^{\mathrm{dl}}}
\right).
\label{eq:downlink_rate}
\end{equation}
We determine the communication costs from the link
conditions at offloading admission and retain them for
the task, even if its remote workload is subsequently
migrated. The reference downlink rate supports this
cost evaluation; downlink resource scheduling is outside
the decision space.

\subsection{Task and Computation Model}
\label{subsec:task_computation_model}

Each user generates computation tasks stochastically.
We characterize the $n$th task of user $i$ by
\begin{equation}
\mathcal{J}_{i,n}=
\left(D_{i,n},C_{i,n},\tau_{i,n},a_{i,n}\right),
\label{eq:task_definition}
\end{equation}
where $D_{i,n}$ and $C_{i,n}$ specify the input size
in bits and computation demand in CPU cycles, respectively.
The task is generated at $a_{i,n}$ and must meet a relative
deadline of $\tau_{i,n}$ seconds. Before dispatch, it
waits in the user's FIFO task queue. If dispatch occurs
in slot $t_{i,n}^{\mathrm d}$ at time $u_{i,n}$, its
waiting time is $W^{\mathrm u}_{i,n}=u_{i,n}-a_{i,n}$,
where both timestamps are measured in seconds.

At dispatch, the user assigns a fraction
$\rho_{i,n}\in[0,1]$ of the task to remote execution,
leaving $(1-\rho_{i,n})C_{i,n}$ cycles for its local
processor and $\rho_{i,n}C_{i,n}$ cycles for the MEC.
Here, $\rho_{i,n}$ refers to the fraction actually
offloaded: a request rejected before upload, or dispatched
without a service connection, falls back to local execution
with $\rho_{i,n}=0$ and no upload energy. For an admitted
request, the serving satellite
$s_{i,n}^{\mathrm{up}}=b_i^{+}(t_{i,n}^{\mathrm d})$
receives the remote workload, whose execution server
may later change through handover.

\subsubsection{Local execution}

Since each terminal serves its scheduled local workloads
serially, dispatch does not necessarily coincide with the
start of local computation. Let $A_i$ denote the processor's
next available time at dispatch and $f_i$ its processing
rate in CPU cycles per second. A nonempty local branch
waits for $W^{\mathrm{loc}}_{i,n}=\max\{A_i-u_{i,n},0\}$
before receiving service, giving
\begin{equation}
T^{\mathrm{loc}}_{i,n}=
W^{\mathrm u}_{i,n}
+W^{\mathrm{loc}}_{i,n}
+\frac{(1-\rho_{i,n})C_{i,n}}{f_i}.
\label{eq:local_delay}
\end{equation}
For the same workload, the local computing energy follows
\begin{equation}
E^{\mathrm{loc}}_{i,n}
=\kappa_i(1-\rho_{i,n})C_{i,n}f_i^2,
\label{eq:local_energy}
\end{equation}
where $\kappa_i$ captures the effective switched
capacitance of the processor. In the full-offloading
case, $\rho_{i,n}=1$, no local branch is scheduled and
we set $T^{\mathrm{loc}}_{i,n}=0$.

\subsubsection{Satellite execution}

At satellite $s$, remote workloads share a finite MEC
queue and a limited number of processing slots.
To distinguish resident tasks from those receiving
service, we denote them by $\mathcal{Q}_s(t)$ and
$\mathcal{P}_s(t)\subseteq\mathcal{Q}_s(t)$, respectively.
With total task capacity $Q_s^{\max}$ and concurrency
limit $K_s$, the server satisfies
\begin{equation}
|\mathcal{Q}_s(t)|\leq Q_s^{\max},
\qquad
|\mathcal{P}_s(t)|\leq K_s.
\label{eq:mec_capacity}
\end{equation}
Here, the task capacity includes both waiting and
processing tasks. At each slot boundary, available processing slots
are filled in FCFS order according to task entry
into the current server's queue.
Within each time slot, the total processing rate
$F_s$ is divided equally among active tasks:
\begin{equation}
f^{\mathrm{MEC}}_{s,q}(t)
=\frac{F_s}{|\mathcal{P}_s(t)|},
\qquad q\in\mathcal{P}_s(t).
\label{eq:mec_rate}
\end{equation}
Holding these shares fixed within the slot, we update
the remaining cycles $C_q^{\mathrm{res}}(t)$ of an
active task $q$ according to
\begin{equation}
C_q^{\mathrm{res}}(t+1)=
\max\left\{
C_q^{\mathrm{res}}(t)
-f^{\mathrm{MEC}}_{s,q}(t)\Delta t,0
\right\},
\label{eq:remaining_workload}
\end{equation}
whereas waiting tasks retain their remaining workloads.
Both $F_s$ and $f^{\mathrm{MEC}}_{s,q}(t)$ are measured
in CPU cycles per second.

In addition to MEC service, a remote branch incurs
communication costs for uploading its input and returning
its output. For $\rho_{i,n}>0$, the transmitted volumes
are $\rho_{i,n}D_{i,n}$ and
$\zeta\rho_{i,n}D_{i,n}$, respectively, where $\zeta$
is the output-to-input size ratio. Evaluating the initial
link at $s=s_{i,n}^{\mathrm{up}}$ and
$t=t_{i,n}^{\mathrm d}$ gives
\begin{align}
T^{\mathrm{up}}_{i,n}
&=\frac{\rho_{i,n}D_{i,n}}
{R^{\mathrm{up}}_{i,s}(t)}
+\frac{d_{i,s}(t)}{c},
\label{eq:upload_delay}\\
T^{\mathrm{dl}}_{i,n}
&=\frac{\zeta\rho_{i,n}D_{i,n}}
{R^{\mathrm{dl}}_{s,i}(t)}
+\frac{d_{i,s}(t)}{c}.
\label{eq:download_delay}
\end{align}

For a remote branch that completes, we distinguish
MEC waiting from processor service across all visited
servers. Let $W^{\mathrm{MEC}}_{i,n}$ accumulate the
intervals from queue entry to service start, or to
reassignment if a task leaves before receiving service.
Let $S^{\mathrm{MEC}}_{i,n}$ accumulate the intervals
from service start to completion or reassignment.
Both are measured on the slotted MEC service timeline
and exclude user-queue waiting, communication costs,
and fixed handover overhead. The MEC component is
\begin{equation}
T^{\mathrm{MEC}}_{i,n}
=
W^{\mathrm{MEC}}_{i,n}
+
S^{\mathrm{MEC}}_{i,n}.
\label{eq:mec_residence}
\end{equation}

With $H_{i,n}$ denoting the accumulated fixed
association-change overhead, the remote-branch delay
becomes
\begin{equation}
\begin{aligned}
T^{\mathrm{rem}}_{i,n}
={}&W^{\mathrm u}_{i,n}
+T^{\mathrm{up}}_{i,n}
+W^{\mathrm{MEC}}_{i,n}\\
&+S^{\mathrm{MEC}}_{i,n}
+T^{\mathrm{dl}}_{i,n}
+H_{i,n}.
\end{aligned}
\label{eq:remote_delay}
\end{equation}

The terminal and MEC process their assigned portions
in parallel. For a task whose required branches complete,
the overall delay therefore follows
\begin{equation}
T_{i,n}=
\max\left\{
T^{\mathrm{loc}}_{i,n},
T^{\mathrm{rem}}_{i,n}
\right\}.
\label{eq:task_delay}
\end{equation}
For local-only execution, the remote branch is absent
and we set $T^{\mathrm{rem}}_{i,n}=0$.

For energy evaluation, we focus on the terminal's
local computing and uplink transmission costs.
With power-amplifier efficiency $\eta_i$ and circuit
power $P_i^{\mathrm{cir}}$, an admitted remote branch
consumes
\begin{equation}
E^{\mathrm{up}}_{i,n}=
\left(\frac{P_i}{\eta_i}+P_i^{\mathrm{cir}}\right)
\frac{\rho_{i,n}D_{i,n}}
{R^{\mathrm{up}}_{i,s}(t)},
\label{eq:upload_energy}
\end{equation}
where $s$ and $t$ again refer to the initial admission
snapshot. No upload energy is incurred for a local-only
task. Together with local computing, the terminal-energy
cost is
\begin{equation}
E_{i,n}=E^{\mathrm{loc}}_{i,n}+E^{\mathrm{up}}_{i,n},
\label{eq:task_energy}
\end{equation}
with $E^{\mathrm{up}}_{i,n}=0$ when $\rho_{i,n}=0$.
Satellite computing and downlink energy do not enter
this terminal-energy objective.

We count a task as successful only when its required
workload completes without branch termination and its
evaluated delay satisfies $T_{i,n}\leq\tau_{i,n}$.
Tasks exceeding this limit, including those expiring
in user or MEC queues, contribute to deadline violations.
We record disconnection-induced terminations separately,
while tasks unresolved at the evaluation horizon remain
pending. For an expired or terminated task, we record
its outcome rather than assign a normal completion delay.

In evaluating performance, we normalize success and
deadline-violation counts by the number of generated
tasks and average successful-task delay over successful
tasks only. Energy per successful task divides the
total terminal energy of all dispatched workloads by
the successful-task count; its numerator includes
energy charged to unsuccessful and unresolved tasks.

\subsection{Handover and Service Model}
\label{subsec:handover_model}

To maintain service as the constellation moves, a user
may retain its association or request a handover to
another visible satellite. Such a request is enabled
when the user is disconnected, its current link becomes
invalid, the remaining visibility falls below
$v_{\mathrm{trig}}$, or the serving MEC reaches its
utilization threshold. The target must meet the minimum
elevation requirement, remaining visibility threshold
$v_{\min}$, and reference uplink SNR threshold
$\Gamma_{\min}$, with the reference SNR evaluated
over total uplink bandwidth $B_s$.

Besides link feasibility, the target must accommodate
the user's unfinished remote workloads. Specifically,
if $m_i(t)$ tasks require migration from the current
MEC, a target satellite $s$ must satisfy
\begin{equation}
|\mathcal{Q}_s(t)|+m_i(t)\leq Q_s^{\max}.
\label{eq:migration_capacity}
\end{equation}
The action mask filters candidate targets using these
conditions and always retains the stay action.
Since other users' association decisions may change
the target queue, we recheck capacity before committing
the handover.

Once committed, the handover updates $b_i^{+}(t)$
and reassigns the user's pending MEC tasks together.
Each task preserves its remaining cycles, generation
time, and deadline, but joins the tail of the target
queue and waits for a processing slot under the same
FCFS rule. If a request fails, the user keeps its old
association as long as that link remains valid.
Otherwise, it becomes disconnected and its stranded
remote tasks are terminated. Local execution remains
available for new tasks, and a feasible association
can restore remote service through reconnection.
We count reconnections separately from handovers.

Each committed handover or reconnection contributes
a fixed service-interruption cost $\delta^{\mathrm{ho}}$.
For a remote task $\mathcal{J}_{i,n}$, let $h_{i,n}$
count the applicable association changes at dispatch
and during subsequent migrations. Its accumulated
overhead is $H_{i,n}=h_{i,n}\delta^{\mathrm{ho}}$.
Each event contributes once to \eqref{eq:remote_delay};
the fixed overhead does not advance the MEC service
timeline or enter \eqref{eq:mec_residence}.
The local branch retains its own processor timeline.

\section{Problem Formulation}
\label{sec:problem_formulation}

Our objective is to jointly determine satellite associations
and partial offloading decisions that improve long-term task
service under time-varying connectivity and shared MEC
resources. Let $k_i(t)$ denote the satellite action of user
$i$, where $k_i(t)=0$ represents staying and
$k_i(t)=s$ requests association with satellite $s$.
The corresponding requested offloading fraction is
$\widehat{\rho}_i(t)$.
The candidate set
$\mathcal{C}_i(t)\subseteq
\mathcal{V}_i(t)\setminus\{b_i^{-}(t)\}$
contains visible alternatives to the current satellite.
The resulting association $b_i^{+}(t)$ and executed task
fraction $\rho_{i,n}$ follow the migration, admission,
and local-execution rules in Section~\ref{sec:system_model}.

To account for both task completion and execution quality,
let $\mathcal{Z}_i^{\mathrm{s}}(t)$ and
$\mathcal{Z}_i^{\mathrm{f}}(t)$ collect the tasks of user $i$
whose outcomes are resolved as successful and failed,
respectively, in slot $t$.
For each successful task, define the normalized cost
\begin{equation}
c_{i,n}
=
\omega_{\mathrm D}\frac{T_{i,n}}{\tau_{i,n}}
+
\omega_{\mathrm E}
\frac{E_{i,n}}{E_{i,n}+E_{\mathrm{ref}}},
\label{eq:normalized_task_cost}
\end{equation}
where $E_{\mathrm{ref}}>0$ is an energy reference,
$\omega_{\mathrm D},\omega_{\mathrm E}\geq0$, and
$\omega_{\mathrm D}+\omega_{\mathrm E}=1$.
The slot utility is then
\begin{equation}
\begin{aligned}
U(t)
={}&\frac{1}{N}\sum_{i\in\mathcal{U}}
\Bigg[
\sum_{n\in\mathcal{Z}_i^{\mathrm{s}}(t)}
(1-c_{i,n})
-
|\mathcal{Z}_i^{\mathrm{f}}(t)|\\
&\hspace{12mm}
-\omega_{\mathrm I}I_i(t)
-\omega_{\mathrm F}F_i(t)
\Bigg]
+\omega_{\mathrm L}J(t),
\end{aligned}
\label{eq:network_service_utility}
\end{equation}
where $I_i(t)\in[0,1]$ is the service-interruption duration
normalized by $\Delta t$, and $F_i(t)$ indicates a failed
association request.
The shared term $J(t)$ is Jain's fairness index of the
reachable MEC loads, with each load averaging normalized
queue occupancy and CPU utilization; $J(t)=0$ when all
reachable servers are idle.
The nonnegative weights
$\omega_{\mathrm I}$, $\omega_{\mathrm F}$, and
$\omega_{\mathrm L}$ specify the service-continuity and
load-balancing preferences.

Let $\pi$ be a causal policy for the joint requests
$\{k_i(t),\widehat{\rho}_i(t)\}_{i\in\mathcal{U}}$.
Furthermore, let $\chi_i(t)$ indicate whether the
association-change conditions in
Section~\ref{subsec:handover_model} hold, and let
$\mathcal{Q}_s^{+}(t)$ denote the MEC queue after migration
and new-task admission.
The joint optimization problem is
\begin{equation}
\begin{aligned}
\mathrm{(P1)}:\quad
\max_{\pi}\quad&
\mathbb{E}_{\pi}
\left[\sum_{t=0}^{\infty}\gamma^t U(t)\right]\\
\mathrm{s.t.}\quad
&\mathrm{C}_1:\;
k_i(t)\in\{0\}\cup\mathcal{C}_i(t),\\
&\mathrm{C}_2:\;
0\leq\widehat{\rho}_i(t)\leq1,\\
&\mathrm{C}_3:\;
k_i(t)=0,\quad\text{if }\chi_i(t)=0,\\
&\mathrm{C}_4:\;
v_{i,s}(t)\geq v_{\min},\
\Gamma^{\mathrm{ref}}_{i,s}(t)\geq\Gamma_{\min},\\
&\mathrm{C}_5:\;
|\mathcal{Q}_s(t)|+m_i(t)\leq Q_s^{\max},\\
&\mathrm{C}_6:\;
|\mathcal{Q}_s^{+}(t)|\leq Q_s^{\max},\
|\mathcal{P}_s(t)|\leq K_s,
\end{aligned}
\label{eq:joint_optimization_problem}
\end{equation}
where $\gamma\in(0,1)$ is the discount factor and
$\Gamma^{\mathrm{ref}}_{i,s}(t)
=P_i g_{i,s}(t)/(N_0B_s)$.
Constraints apply at every slot, with
$\mathrm{C}_4$ and $\mathrm{C}_5$ evaluated for each
requested target $s=k_i(t)\neq0$.

Specifically, $\mathrm{C}_1$ permits staying or selecting
one visible candidate, while $\mathrm{C}_2$ bounds the
requested workload fraction.
Constraint $\mathrm{C}_3$ enforces the association-change
trigger, and $\mathrm{C}_4$ imposes the target-link
requirements.
Constraint $\mathrm{C}_5$ screens migration requests
using the current queue state; capacity is rechecked
at commitment.
Finally, $\mathrm{C}_6$ enforces the queue and processing
limits for executed decisions.
Task deadlines determine the success and failure sets
in \eqref{eq:network_service_utility}.

Problem~\eqref{eq:joint_optimization_problem} couples
discrete satellite selection with continuous workload
allocation. A selected satellite determines the link
and computing conditions, while the offloading fraction
changes the contention experienced by other users.
Moreover, residual workloads and handover-induced
reassignment carry these effects across slots.
Consequently, effective decisions require both a
representation of multi-user contention and an
evaluation of the long-term value of each
satellite--fraction pair. These requirements motivate
the HAN--PDQN framework developed next.

\section{Proposed HAN--PDQN Framework}
\label{sec:proposed_framework}

\subsection{Framework Overview}
\label{subsec:framework_overview}

To address \eqref{eq:joint_optimization_problem}, we develop
HAN--PDQN, which connects a representation of multi-user
contention to parameterized satellite and offloading decisions.
As illustrated in Fig.~\ref{fig:han_pdqn_framework}, the
controller constructs a heterogeneous graph from user demand,
satellite resources, and their serving and visibility relations.
HAN aggregates these measurements into a user representation
that distinguishes established resource sharing from potential
competition through overlapping satellite alternatives.

The representation is combined with the user's local task
and candidate-satellite observations. Conditioned on this input,
PDQN generates an offloading fraction for each satellite
action and evaluates the resulting satellite--fraction pairs.
All users share the decision-network parameters, while their
observations and feasible candidates produce different actions.
The joint requests are then executed under the service rules
in Section~\ref{sec:system_model}, yielding task outcomes and
the next network snapshot. Sections~\ref{subsec:han_contention}
and~\ref{subsec:pdqn_joint_control} develop the representation
and decision mechanisms, respectively;
Section~\ref{subsec:training_execution} describes training
and online execution.

\begin{figure*}[t]
    \centering
    \includegraphics[
        trim=3mm 1mm 0.2mm 2mm,
        clip,
        width=\textwidth
    ]{算法框架图10_23.pdf}
    \caption{HAN--PDQN framework for joint handover and partial
    task offloading.}
    \label{fig:han_pdqn_framework}
\end{figure*}

\subsection{Heterogeneous Graph-Based Contention Representation}
\label{subsec:han_contention}

At slot $t$, we construct
$\mathcal{G}_t=(\mathcal{V},\mathcal{E}_t,\mathbf{X}_t,
\mathbf{E}_t)$, where $\mathcal{V}$ contains user and satellite
nodes, $\mathcal{E}_t$ contains typed edges, and
$\mathbf{X}_t$ and $\mathbf{E}_t$ collect node and edge
features. User features describe location, connection status,
task demand, urgency, and backlog. Satellite features describe
location, MEC utilization, queue occupancy, associated users,
and available computing resources. Visibility edges carry
distance, elevation, SNR, reference rate, and remaining
visibility; serving edges identify existing associations
and their service duration. Nearby-user and ISL edges provide
geographical and constellation-neighborhood information.
User--satellite relations are represented in both directions.

The five meta-paths in Table~\ref{tab:han_metapaths}
organize these relations into complementary service contexts.
The serving-based U--S--U path aggregates users attached to
the same MEC, exposing established contention. The
visibility-based U--S--U path connects users with common
satellite alternatives, exposing demand that may converge
on the same server after subsequent decisions. This
distinction retains both the current sharing structure and
the opportunities through which that structure can change.
The remaining paths incorporate neighboring demand,
overlapping satellite coverage, and ISL neighborhoods.

\begin{table}[H]
    \caption{Meta-Paths Used for Contention Representation}
    \label{tab:han_metapaths}
    \centering
    \footnotesize
    \renewcommand{\arraystretch}{1.13}
    \begin{tabular}{p{0.38\columnwidth}p{0.53\columnwidth}}
        \hline
        Meta-path and relation & Service context \\
        \hline
        U--S--U (serving) &
        Users sharing a serving MEC \\
        U--S--U (visibility) &
        Users with common satellite alternatives \\
        U--U (proximity) &
        Demand from neighboring users \\
        S--U--S (visibility) &
        Satellites visible to common users \\
        S--S--S (ISL) &
        Satellite context across two ISL hops \\
        \hline
    \end{tabular}
\end{table}

Following the hierarchical attention principle
in~\cite{wang2019han}, we first aggregate neighbors along
each meta-path and then combine the resulting semantic
representations. At each HAN layer, type-specific projections
map node features into a common hidden space. Each meta-path
starts from these projected features and propagates messages
successively along its constituent relations.

For one relation $r$ at a given meta-path hop, let
$\mathcal{N}_r(v)$ denote the incoming neighbors of node $v$.
For attention head $h$, we form the relation input
\begin{equation}
\mathbf{c}_{uv}^{r,h}
=
\left[
\mathbf{W}_{\mathrm s}^{r,h}\mathbf{h}_u
\Vert
\mathbf{W}_{\mathrm d}^{r,h}\mathbf{h}_v
\Vert
\mathbf{W}_{\mathrm e}^{r,h}\mathbf{e}_{uv}^{r}
\right],
\label{eq:han_attention_input}
\end{equation}
where $\Vert$ denotes concatenation and
$\mathbf{e}_{uv}^{r}$ is the edge feature.
The attention score and normalized weight are
\begin{equation}
\begin{aligned}
e_{uv}^{r,h}
&=
\operatorname{LeakyReLU}
\left(
(\mathbf{a}^{r,h})^{\mathsf T}
\mathbf{c}_{uv}^{r,h}
\right),\\
\alpha_{uv}^{r,h}
&=
\frac{\exp(e_{uv}^{r,h})}
{\sum_{w\in\mathcal{N}_r(v)}
 \exp(e_{wv}^{r,h})}.
\end{aligned}
\label{eq:han_relation_attention}
\end{equation}
The matrices and attention vectors are specific to the layer,
meta-path, and hop; these indices are suppressed for clarity.
The $H$ heads produce the message and residual update
\begin{equation}
\begin{aligned}
\mathbf{m}_v^r
&=
\mathop{\Vert}_{h=1}^{H}
\sum_{u\in\mathcal{N}_r(v)}
\alpha_{uv}^{r,h}
\mathbf{W}_{\mathrm s}^{r,h}\mathbf{h}_u,\\
\mathbf{h}_v^{+}
&=
\operatorname{LN}
\left(
\mathbf{h}_v+
\operatorname{ELU}(\mathbf{m}_v^r)
\right).
\end{aligned}
\label{eq:han_hop_update}
\end{equation}
Applying \eqref{eq:han_hop_update} along a meta-path $\phi$
yields its output $\mathbf{h}_v^\phi$.
Thus, on a U--S--U path, user demand is first aggregated
at the intermediate satellite and then propagated to
the destination users.

Let $\Phi_\nu$ contain the meta-paths ending at node type
$\nu\in\{\mathrm U,\mathrm S\}$, and let $\mathcal{V}_\nu$
be the corresponding node set. Semantic attention assigns
each applicable path the score
\begin{equation}
w_{\phi,\nu}
=
\frac{1}{|\mathcal{V}_\nu|}
\sum_{v\in\mathcal{V}_\nu}
\mathbf{q}_\nu^{\mathsf T}
\tanh\left(
\mathbf{W}_\nu\mathbf{h}_v^\phi+\mathbf{b}_\nu
\right).
\label{eq:han_semantic_score}
\end{equation}
The path weights and fused node representation are
\begin{equation}
\begin{aligned}
\beta_{\phi,\nu}
&=
\frac{\exp(w_{\phi,\nu})}
{\sum_{\phi'\in\Phi_\nu}\exp(w_{\phi',\nu})},\\
\mathbf{z}_v
&=
\mathbf{W}^{\mathrm{out}}_\nu
\sum_{\phi\in\Phi_\nu}
\beta_{\phi,\nu}\mathbf{h}_v^\phi
+\mathbf{b}^{\mathrm{out}}_\nu.
\end{aligned}
\label{eq:han_semantic_fusion}
\end{equation}
These semantic weights are shared by nodes of the same type
within a graph snapshot, while node-level weights depend
on the contributing neighbors.

Stacking HAN layers allows satellite representations formed
through S--U--S and S--S--S paths to enter user-directed
propagation in the next layer. Consequently, the final
user embedding $\mathbf{z}_i(t)$ incorporates the user's
own demand, competing users, and the surrounding satellite
resource context. This embedding supplies the common
relational input to both PDQN components.

\subsection{Reinforcement Learning Formulation and PDQN-Based Control}
\label{subsec:pdqn_joint_control}

We model the underlying sequential control problem as
a Markov decision process
$\mathcal{M}=(\mathcal{S},\mathcal{A},\mathcal{P},R,\gamma)$,
where $\mathcal{S}$ is the network state space,
$\mathcal{A}$ is the joint action space,
$\mathcal{P}$ describes the state transition law,
$R$ is the network utility in
\eqref{eq:network_service_utility}, and $\gamma$ is
the discount factor.
The controller forms a graph-augmented observation
for each user and applies shared PDQN networks
to generate the joint requests.
The state, action, and reward are specified below.

\subsubsection{State and Observation}

The network state at the beginning of slot $t$
is represented by
\begin{equation}
\mathbf{s}(t)
=
\big(
\boldsymbol{\zeta}(t),
\mathbf{b}^{-}(t),
\mathcal{W}(t),
\boldsymbol{\ell}(t),
\boldsymbol{h}(t)
\big),
\label{eq:rl_network_state}
\end{equation}
where $\boldsymbol{\zeta}(t)$ collects the orbital,
user-location, and channel conditions;
$\mathbf{b}^{-}(t)$ contains the current satellite
associations; and $\mathcal{W}(t)$ records the
ordered user and MEC queues, active tasks,
residual workloads, and task timing attributes.
Furthermore, $\boldsymbol{\ell}(t)$ records local
processor availability, while $\boldsymbol{h}(t)$
contains connection states and the timers governing
association changes and service interruption.
Together, these quantities determine the service
evolution under the model in
Section~\ref{sec:system_model}.

For decision making, user $i$ is represented by
a raw observation $\mathbf{o}_i(t)$ containing
its location and connection status, serving-link
measurements, head-of-line task attributes,
task urgency, queue backlog, and ordered
candidate-satellite features.
Each candidate entry includes its satellite
identifier, distance, elevation, SNR, remaining
visibility, and MEC utilization.
Unused candidate entries are zero-padded.

The observation is combined with the user embedding
from Section~\ref{subsec:han_contention}:
\begin{equation}
\widetilde{\mathbf{o}}_i(t)
=
\left[
\mathbf{o}_i(t)\Vert\mathbf{z}_i(t)
\right].
\label{eq:fused_pdqn_observation}
\end{equation}
Here, $\mathbf{o}_i(t)$ preserves the task and
candidate-specific measurements, whereas
$\mathbf{z}_i(t)$ summarizes the surrounding
resource-sharing relationships.
The resulting vector is the input to both
PDQN components.

\subsubsection{Action Space}

Each user jointly selects a satellite action
and its associated offloading fraction.
Let $K$ be the maximum number of candidate slots,
and let $s_{i,j}(t)$ denote the satellite occupying
slot $j$.
The discrete index $d_i(t)\in\{0,\ldots,K\}$
reserves index $0$ for staying with the current
association.

Let $m_{i,j}(t)\in\{0,1\}$ indicate whether
index $j$ satisfies the current action conditions.
The admissible index set is
$\mathcal{K}_i(t)=\{j:m_{i,j}(t)=1\}$,
which always contains the stay index.
The parameterized action is
\begin{equation}
\begin{aligned}
a_i(t)
&=
\big(d_i(t),\widehat{\rho}_i(t)\big)
\in\mathcal{A}_i(t),\\
\mathcal{A}_i(t)
&=
\bigcup_{j\in\mathcal{K}_i(t)}
\left(\{j\}\times[0,1]\right).
\end{aligned}
\label{eq:rl_parameterized_action_space}
\end{equation}
Its discrete component maps to the satellite
decision in \eqref{eq:joint_optimization_problem} as
\[
k_i(t)
=
\begin{cases}
0, & d_i(t)=0,\\
s_{i,d_i(t)}(t), & d_i(t)>0.
\end{cases}
\]
Thus, the continuous component specifies how much
work to offload through the association resulting
from the discrete request.
In particular, index $0$ can still offload through
the current serving satellite.
Requests below $\rho_{\min}$ invoke local execution.

The candidate features, parameter branches, and
action mask use the same satellite-to-index mapping
within each slot.
Unused slots are masked, together with candidates
that fail the trigger, link, or migration checks
in $\mathrm{C}_1$ and
$\mathrm{C}_3$--$\mathrm{C}_5$.
Shared capacity is checked again when requests
are committed, enforcing $\mathrm{C}_6$.

The joint action is
$\mathbf{a}(t)=(a_i(t))_{i\in\mathcal{U}}$.
Following its execution, orbital motion,
stochastic task arrivals, workload processing,
and association changes produce the next state
according to
$\mathbf{s}(t+1)\sim
\mathcal{P}(\cdot\mid\mathbf{s}(t),\mathbf{a}(t))$.
The next observations and action masks are then
constructed from the updated network conditions.

\subsubsection{Reward Function}

The reward combines task-service quality with
the operating costs of maintaining satellite access.
Using the normalized task cost
$c_{i,n}$ in \eqref{eq:normalized_task_cost},
the feedback assigned to user $i$ is
\begin{equation}
\begin{aligned}
r_i(t)
={}&
\sum_{n\in\mathcal{Z}_i^{\mathrm{s}}(t)}
(1-c_{i,n})
-
|\mathcal{Z}_i^{\mathrm{f}}(t)|\\
&-\omega_{\mathrm I}I_i(t)
-\omega_{\mathrm F}F_i(t)
+\omega_{\mathrm L}J(t).
\end{aligned}
\label{eq:pdqn_user_reward}
\end{equation}
Successful tasks contribute a positive reward
discounted by their normalized delay and
terminal energy cost, whereas each failed task
incurs a unit penalty.
The next two terms penalize service interruption
and failed association requests.
The shared term $J(t)$ provides feedback on
the balance of reachable MEC loads.

The user rewards satisfy
\[
\frac{1}{N}
\sum_{i\in\mathcal{U}}r_i(t)=U(t),
\]
linking the learning feedback to the network
objective in
\eqref{eq:joint_optimization_problem}.
The shared Q-network uses these user-level returns
to approximate the value of each user's action,
with the HAN representation providing its
multi-user resource context.
Task rewards are assigned when outcomes are resolved,
so feedback in slot $t$ can include workloads
dispatched in earlier slots.
The settlement timing is specified in
Section~\ref{subsec:training_execution}.

\subsubsection{PDQN-Based Decision}

The action space combines a discrete satellite
index with a continuous parameter associated
with that index.
Following the parameterized action formulation
in~\cite{xiong2018pdqn}, we use a parameter
network to generate one requested fraction
for each discrete branch:
\begin{equation}
\widehat{\rho}_{i,j}(t)
=
\mu_{\psi,j}
\left(\widetilde{\mathbf{o}}_i(t)\right)
=
\sigma\left(
f_{\psi,j}
\left(\widetilde{\mathbf{o}}_i(t)\right)
\right),
\label{eq:pdqn_parameter_generation}
\end{equation}
where $\sigma$ is the sigmoid function and
$\psi$ collects the parameter-network weights.
The output mapping satisfies the fraction bound
in $\mathrm{C}_2$.

Given these fractions, the Q-network evaluates
each candidate index--fraction pair:
\begin{equation}
q_{i,j}(t)
=
Q_{\theta}
\left(
\widetilde{\mathbf{o}}_i(t),
j,
\widehat{\rho}_{i,j}(t)
\right).
\label{eq:pdqn_pair_value}
\end{equation}
Its inputs are the fused observation,
a one-hot encoding of $j$, and the scalar
fraction generated for that branch.
Its output estimates the discounted user return
obtained by selecting this pair and following
the learned policy thereafter.
The parameters $\psi$ and $\theta$ are shared
across users.

During online execution, the controller selects
the admissible pair with the largest estimated value:
\begin{equation}
\begin{aligned}
d_i^{\star}(t)
&=
\arg\max_{j\in\mathcal{K}_i(t)}
q_{i,j}(t),\\
\widehat{\rho}_i(t)
&=
\widehat{\rho}_{i,d_i^{\star}(t)}(t).
\end{aligned}
\label{eq:pdqn_joint_selection}
\end{equation}
Consequently, satellite selection accounts for
the workload fraction proposed for each candidate.
The same contention representation conditions
both fraction generation and pair evaluation,
connecting the amount of offloaded work to
the communication and computing conditions
of its destination.
The replay-based updates and exploration strategy
are presented next.

\subsection{Training and Online Execution}
\label{subsec:training_execution}

Training uses an experience replay buffer $\mathcal{D}$
of joint transitions, each containing the user tuples
$(\widetilde{\mathbf{o}},d,\widehat{\rho},r,
\widetilde{\mathbf{o}}',\delta,\mathbf{m},\mathbf{m}')$.
Here, $\widehat{\rho}$ is the requested fraction,
the primed quantities describe the next observation,
and $\delta$ indicates a true terminal transition.
The buffer stores the encoded observations used for
the decisions. A minibatch of $B$ joint transitions
provides $BN$ user samples for the shared PDQN networks.

Task feedback is assigned when an outcome is resolved.
For local-only execution, the reserved completion time
determines the outcome at dispatch; remote outcomes
are evaluated upon MEC completion, expiration, or
termination. Interruption, association-failure, and
load-balance feedback are assigned in the corresponding
slot. Thus, the replay reward follows
\eqref{eq:pdqn_user_reward}, including outcomes from
previously dispatched workloads.

Let $Q_{\bar{\theta}}$ and $\mu_{\bar{\psi},j}$
denote the target networks. For a sampled next observation,
the online networks select
\begin{equation}
d'
=
\arg\max_{j:m'_j=1}
Q_\theta\left(
\widetilde{\mathbf{o}}',j,
\mu_{\psi,j}(\widetilde{\mathbf{o}}')
\right).
\label{eq:pdqn_next_action}
\end{equation}
The target networks then evaluate this choice:
\begin{equation}
y
=
r+\gamma(1-\delta)
Q_{\bar{\theta}}\left(
\widetilde{\mathbf{o}}',d',
\mu_{\bar{\psi},d'}(\widetilde{\mathbf{o}}')
\right).
\label{eq:pdqn_td_target}
\end{equation}
The Q-network minimizes
\begin{equation}
\mathcal{L}_Q(\theta)
=
\mathbb{E}_{\mathcal{D}}
\left[
\left(
Q_\theta(\widetilde{\mathbf{o}},d,\widehat{\rho})-y
\right)^2
\right].
\label{eq:pdqn_q_loss}
\end{equation}
Time-limit truncation retains bootstrapping, since
unfinished workloads continue to have future value.

The parameter branches are updated through the Q-network
using the objective
\begin{equation}
\mathcal{L}_{\mu}(\psi)
=
-\eta_{\mu}
\mathbb{E}_{\mathcal{D}}
\left[
\max_{j:m_j=1}
Q_\theta\left(
\widetilde{\mathbf{o}},j,
\mu_{\psi,j}(\widetilde{\mathbf{o}})
\right)
\right],
\label{eq:pdqn_parameter_loss}
\end{equation}
where $\eta_{\mu}>0$ scales the parameter objective.
During this update, $\theta$ is held constant and
gradients propagate through the continuous parameter
of the maximizing feasible branch.
Both optimizers use gradient clipping.
Every $C_{\mathrm{tar}}$ learning updates, the target
parameters are copied from the online networks:
$\bar{\theta}\leftarrow\theta$ and
$\bar{\psi}\leftarrow\psi$.

Exploration samples a feasible discrete index and a
continuous fraction uniformly from $[0,1]$.
After the initial replay warm-up, this random pair is
selected with probability $\epsilon$; otherwise, the
controller uses \eqref{eq:pdqn_joint_selection}.
The exploration probability decreases during training.
Algorithm~\ref{alg:han_pdqn_training} summarizes the procedure.

\begin{algorithm}[t]
\caption{HAN--PDQN Training}
\label{alg:han_pdqn_training}
\begin{algorithmic}[1]
\Require HAN encoder, candidate limit $K$, replay capacity,
batch size $B$, target period $C_{\mathrm{tar}}$
\State Initialize $\theta,\psi$ and replay buffer $\mathcal{D}$
\State Set $\bar{\theta}\gets\theta$, $\bar{\psi}\gets\psi$
\For{each training episode}
    \State Reset the environment and construct $\mathcal{G}_0$
    \State Form $\widetilde{\mathbf{o}}_i(0)$ and
    $\mathbf{m}_i(0)$ for all users
    \For{each decision slot $t$}
        \For{each user $i\in\mathcal{U}$}
            \State Select a feasible pair by exploration or
            \eqref{eq:pdqn_joint_selection}
        \EndFor
        \State Execute the joint requests under the service model
        \State Obtain $r_i(t)$, $\widetilde{\mathbf{o}}_i(t+1)$,
        and $\mathbf{m}_i(t+1)$
        \State Store the joint transition in $\mathcal{D}$
        \If{replay warm-up is complete}
            \State Sample a minibatch from $\mathcal{D}$
            \State Compute targets using
            \eqref{eq:pdqn_next_action}--\eqref{eq:pdqn_td_target}
            \State Update $\theta$ using \eqref{eq:pdqn_q_loss}
            \State Update $\psi$ using \eqref{eq:pdqn_parameter_loss}
            \State Copy online parameters to targets every
            $C_{\mathrm{tar}}$ updates
        \EndIf
        \If{episode terminates or reaches its time limit}
            \State \textbf{break}
        \EndIf
    \EndFor
\EndFor
\end{algorithmic}
\end{algorithm}

For online operation, the controller constructs the graph,
computes the fused user observations, and selects the
masked maximizing pairs with $\epsilon=0$.
Association changes are resolved before new upload
requests are admitted and allocated bandwidth, following
Section~\ref{sec:system_model}.
This sequence preserves the candidate mapping used
at decision time while applying capacity checks to
the committed actions.

For computational cost, let $L$ be the number of HAN
layers and $D_{\mathrm h}$ the hidden width.
Let $V_{\mathrm{mp}}$ and $E_{\mathrm{mp}}$ count node
and edge visits across the meta-path hops in one layer.
For a fixed number of heads, graph encoding requires
$O(L(V_{\mathrm{mp}}D_{\mathrm h}^2+
E_{\mathrm{mp}}D_{\mathrm h}))$ operations.
If $P_\mu$ is the total parameter-branch weight count
and $P_Q$ is the Q-network weight count, evaluating all
users adds $O(N[P_\mu+(K+1)P_Q])$.
A replay update over $B$ joint transitions has network cost
$O(BN[P_\mu+(K+1)P_Q])$.
Graph construction contributes a separate cost determined
by visibility and neighborhood generation.
Online execution uses graph construction, encoding,
and forward action evaluation; replay updates contribute
to training cost.

\FloatBarrier
\LEOPageTopLayout
\section{Performance Evaluation}
\label{sec:performance_evaluation}

\subsection{Experimental Setup and Evaluation Protocol}
\label{subsec:experimental_setup}

We evaluate HAN--PDQN in a multi-user LEO satellite edge
computing network comprising 66 satellites distributed
across six orbital planes. The orbital altitude and
inclination are 550~km and $53^\circ$, respectively.
Ground users are distributed around
$(39.9^\circ\mathrm{N},116.4^\circ\mathrm{E})$, and the
population varies from 20 to 40 users.
User locations and orbital starting times are resampled
at episode initialization. Each episode contains
512 slots, with a slot duration of 1~s.
Tables~\ref{tab:simulation_parameters} and~\ref{tab:training_parameters}
summarize the system and training parameters, respectively.

% Place the two parameter tables together, outside the text columns.
\begin{table*}[t]
\centering
\noindent
\begin{minipage}[t]{0.485\textwidth}
\vspace{0pt}
\LEOTableSetup
\caption{System Parameters}
    \label{tab:simulation_parameters}
    \centering
    \footnotesize
    \renewcommand{\arraystretch}{1.08}
    \setlength{\tabcolsep}{4pt}
    \begin{tabular}{p{0.60\linewidth}p{0.32\linewidth}}
        \hline
        Parameter & Value \\
        \hline
        Orbital planes / satellites per plane & $6/11$ \\
        Orbital altitude / inclination & $550$ km / $53^\circ$ \\
        Number of users, $N$ & $20,25,30,35,40$ \\
        User deployment radius & $5^\circ$ \\
        Slot duration, $\Delta t$ & $1$ s \\
        Episode length & $512$ slots \\
        Carrier frequency, $f_c$ & $20$ GHz \\
        Uplink bandwidth per satellite, $B_s$ & $10$ MHz \\
        User / satellite transmit power & $24/40$ dBm \\
        User / satellite antenna gain & $38.5/34$ dBi \\
        Noise temperature / noise figure & $354.81$ K / $2$ dB \\
        Local processing rate, $f_i$ & $1$ Gcycle/s \\
        Total MEC processing rate, $F_s$ & $10$ Gcycle/s \\
        MEC capacity / concurrency limit & $6/2$ tasks \\
        Task arrival probability & $0.35$ per user-slot \\
        Minimum elevation, $\theta_{\min}$ & $10^\circ$ \\
        Visibility thresholds, $v_{\min},v_{\mathrm{trig}}$
        & $60$ s \\
        Target SNR threshold, $\Gamma_{\min}$ & $-5$ dB \\
        Association-change overhead, $\delta^{\mathrm{ho}}$
        & $0.6$ s \\
        Output-to-input size ratio, $\zeta$ & $0.1$ \\
        Amplifier efficiency / circuit power & $0.38$ / $0.05$ W \\
        Local energy coefficient, $\kappa_i$ & $10^{-27}$ \\
        Minimum effective offloading fraction & $0.05$ \\
        \hline
    \end{tabular}
\end{minipage}\hfill%
\begin{minipage}[t]{0.485\textwidth}
\vspace{0pt}
\LEOTableSetup
\caption{Algorithm and Training Parameters}
    \label{tab:training_parameters}
    \centering
    \footnotesize
    \renewcommand{\arraystretch}{1.08}
    \setlength{\tabcolsep}{4pt}
    \begin{tabular}{p{0.60\linewidth}p{0.32\linewidth}}
        \hline
        Parameter & Value \\
        \hline
        Interaction budget per run & $150{,}000$ steps \\
        Training seeds & $42$--$46$ \\
        Evaluation episodes per seed & $5$ \\
        Discount factor, $\gamma$ & $0.99$ \\
        HAN layers / attention heads & $2/4$ \\
        HAN hidden / output dimension & $64/64$ \\
        PDQN Q-network hidden widths & $256,128$ \\
        PDQN parameter-branch hidden widths & $128,64$ \\
        PDQN learning rate & $10^{-3}$ \\
        HAN--PDQN / PDQN replay batch size & $512/128$ \\
        PDQN replay capacity & $50{,}000$ \\
        PDQN replay warm-up & $1{,}000$ steps \\
        PDQN target-copy interval & $500$ updates \\
        PDQN exploration probability & $1.0\rightarrow0.02$ \\
        MAPPO learning rate & $10^{-4}$ \\
        MAPPO rollout length / update epochs & $1024/4$ \\
        MAPPO minibatch size & $512$ \\
        MADDPG actor / critic learning rate
        & $10^{-4}/(3\times10^{-4})$ \\
        MADDPG replay batch size & $128$ \\
        MADDPG target soft-update coefficient & $0.005$ \\
        \hline
    \end{tabular}
\end{minipage}
\end{table*}

Each user generates a task independently with probability
0.35 per slot. Tasks are classified as light, medium,
and heavy with probabilities 0.3, 0.5, and 0.2.
For light tasks, the input size, computation demand,
and deadline are uniformly sampled from
$[1,3]$~Mbit, $[0.3,0.8]$~Gcycle, and $[1.5,3]$~s,
respectively. The corresponding ranges are
$[5,10]$~Mbit, $[1.2,2.5]$~Gcycle, and $[2.5,5.5]$~s
for medium tasks, and
$[12,25]$~Mbit, $[3,7]$~Gcycle, and $[4,8]$~s
for heavy tasks.
Each MEC provides a total processing rate of
10~Gcycle/s through two 5-GHz cores.
The utility weights are
$(\omega_{\mathrm D},\omega_{\mathrm E},
\omega_{\mathrm I},\omega_{\mathrm F},\omega_{\mathrm L})
=(0.60,0.40,0.30,0.20,0.05)$,
with $E_{\mathrm{ref}}=1$~J.

We compare HAN--PDQN with six benchmarks:
\begin{itemize}
    \item \textit{HAN+MAPPO}: A graph-assisted multi-agent
    proximal policy optimization controller for joint
    satellite selection and offloading.

    \item \textit{MAPPO}: A multi-agent proximal policy
    optimization controller using raw user, task, and
    candidate-satellite observations.

    \item \textit{MADDPG}: A multi-agent deep deterministic
    policy gradient controller using raw observations
    to generate satellite and offloading decisions.

    \item \textit{PDQN}: A parameterized Q-learning
    controller using raw observations without HAN
    representations.

    \item \textit{Min-Distance}: A distance-based policy
    that requests the nearest admissible satellite
    under the handover conditions and requests full
    offloading. It retains the current association
    when no shorter admissible alternative is available.

    \item \textit{Full-Local}: A policy that processes
    all tasks locally and issues no handover requests.
\end{itemize}
All methods operate under the same task, communication,
MEC service, and admission models.

For each user population, the learning methods are trained
with seeds 42--46 and an interaction budget of
$1.5\times10^5$ environment steps per run.
Table~\ref{tab:training_parameters} reports the principal
training settings.
HAN--PDQN and PDQN use replay batch sizes of 512 and 128,
respectively. Periodic model evaluation is performed
every 25,000 environment steps for the PDQN and MAPPO
variants, and every ten episodes for MADDPG.
Checkpoints are selected according to their mean
evaluation return.

Each selected policy is evaluated over five episodes
per training seed, using matched scenario seeds across
methods. Performance metrics are first averaged over
these episodes and then summarized across the five
training seeds. Shaded regions and error bars represent
95\% bootstrap confidence intervals at the training-seed
level.

The reported metrics follow
Section~\ref{subsec:task_computation_model}.
Successful-task delay is averaged over successful tasks,
whereas the deadline violation rate is normalized by
the number of generated tasks.
Energy per successful task divides the total accounted
terminal energy by the successful-task count.
We additionally report cumulative blocked user-time,
total terminal energy per episode, and handovers per
user-minute to characterize service continuity and
operating costs.

\begin{figure}[!t]
\LEOFigureSetup
\centering
    \includegraphics[
        trim=2.5mm 5mm 0mm 1mm,
        clip,
        width=\linewidth
    ]{algorithm_convergence_u30.pdf}
    \caption{Training returns with 30 users and 512 steps
    per episode. Curves show five-seed means with 95\%
    bootstrap confidence intervals. No smoothing is applied;
    MAPPO variants retain their two-episode logging intervals.}
    \label{fig:algorithm_convergence}
\end{figure}

\begin{figure}[!t]
\LEOFigureSetup
\centering
    \includegraphics[
        trim=5.5mm 0.8mm 0mm 5mm,
        clip,
        width=\linewidth
    ]{convergence_u30.pdf}
    \caption{Hyperparameter sensitivity with 30 users:
    (a) learning rate with batch size 512;
    (b) replay batch size with learning rate $10^{-3}$.
    Curves show five-seed means and 95\% bootstrap confidence
    intervals after a three-episode trailing mean.}
    \label{fig:hyperparameter_convergence}
\end{figure}

\begin{figure*}[t]
\LEOFigureSetup
\centering
    \includegraphics[
        trim=1.8mm 3.3mm 1.4mm 1.8mm,
        clip,
        width=\linewidth
    ]{figure1_line_metrics.pdf}
    \caption{Task performance versus the number of users:
    (a) average successful-task delay;
    (b) deadline violation rate;
    (c) terminal energy per successful task.
    Curves show five-training-seed means, and shaded regions
    indicate 95\% bootstrap confidence intervals.}
    \label{fig:task_performance}
\end{figure*}

\subsection{Training Convergence and Hyperparameter Sensitivity}
\label{subsec:training_performance}

Fig.~\ref{fig:algorithm_convergence} compares the training
returns of the five learning methods with 30 users.
The episode return accumulates the user-averaged slot
reward over an episode, and each curve averages the
corresponding returns across training seeds.
The figure includes complete episodes within the
interaction budget.

HAN+MAPPO and MAPPO improve rapidly during the initial
training stage and subsequently approach lower return
levels. HAN--PDQN and PDQN exhibit a more pronounced
increase during the intermediate stage, after which
HAN--PDQN maintains the highest mean return.
Averaged over the final 20 complete episodes,
HAN--PDQN achieves a return of 87.78, compared with
76.85 for PDQN. The corresponding returns of HAN+MAPPO
and MAPPO are approximately 59.30 and 57.86.
MADDPG exhibits lower late-stage returns and greater
variation across seeds.
The comparison therefore favors HAN--PDQN in the
quality of the learned policy under the allocated
interaction budget.

We next examine the sensitivity of HAN--PDQN to its
learning rate and replay batch size in
Fig.~\ref{fig:hyperparameter_convergence}.
For the learning-rate experiment, the batch size is
fixed at 512 and the rate is selected from
$\{10^{-4},3\times10^{-4},10^{-3},3\times10^{-3}\}$.
For the batch-size experiment, the learning rate is
fixed at $10^{-3}$ and the batch size is selected from
$\{128,256,512,1024\}$.
Each setting uses five training seeds.

In Fig.~\ref{fig:hyperparameter_convergence}(a), learning
rates between $10^{-4}$ and $10^{-3}$ attain similar
late-stage returns. Increasing the rate to
$3\times10^{-3}$ produces a substantially lower mean
return and a wider confidence interval.
This behavior indicates sensitivity to aggressive
updates when the value estimates and continuous
action parameters evolve together.

Fig.~\ref{fig:hyperparameter_convergence}(b) shows that
batch sizes 128, 512, and 1024 reach comparable return
levels, whereas batch size 256 yields a lower mean
and greater variation across seeds.
The dependence on batch size is therefore nonmonotonic
over the examined settings.
The default combination of learning rate $10^{-3}$
and batch size 512 attains a high return with a
consistent five-seed training profile and is used
for the subsequent HAN--PDQN comparisons.



\begin{figure*}[t]
\LEOFigureSetup
\centering
    \includegraphics[
        trim=3mm 3.3mm 1.4mm 1.8mm,
        clip,
        width=\linewidth
    ]{figure2_bar_metrics.pdf}
    \caption{Service continuity and operating costs:
    (a) blocked user-time per episode;
    (b) total terminal energy per episode;
    (c) handovers per user-minute.
    Bars show five-training-seed means with 95\% bootstrap
    confidence intervals. Full-Local is omitted from (a)
    and performs no handovers in (c).}
    \label{fig:service_operating_costs}
\end{figure*}

\subsection{Task Performance under User Scaling}
\label{subsec:user_scaling_performance}

Fig.~\ref{fig:task_performance} evaluates task service
as the user population increases from 20 to 40.
The per-user arrival process remains unchanged, so
a larger population increases the aggregate demand
for the shared uplink and MEC resources.
HAN--PDQN achieves the lowest mean successful-task
delay, deadline violation rate, and energy per successful
task at every evaluated population.

In Fig.~\ref{fig:task_performance}(a), the mean
successful-task delay of HAN--PDQN increases from
1554.9~ms at 20 users to 1906.7~ms at 40 users.
The corresponding PDQN values are 1608.2 and 1939.9~ms.
At 30 users, HAN--PDQN achieves 1779.7~ms, compared
with 1820.7~ms for PDQN and 2053.8~ms for HAN+MAPPO.
The increase with population is consistent with
greater contention for upload bandwidth and MEC
service. Throughout this range, the proposed controller
maintains a lower delay despite the growing demand.

The reliability difference is more pronounced in
Fig.~\ref{fig:task_performance}(b).
At 30 users, the deadline violation rate is 6.89\%
for HAN--PDQN and 9.47\% for PDQN, a reduction of
2.58 percentage points.
HAN+MAPPO and MAPPO yield 18.38\% and 18.83\%,
respectively, under the same population.
At 40 users, HAN--PDQN maintains a violation rate
of 10.99\%, compared with 13.08\% for PDQN,
24.24\% for HAN+MAPPO, and 37.47\% for Min-Distance.
These results show that the delay advantage is
accompanied by improved deadline compliance.

Fig.~\ref{fig:task_performance}(c) further shows that
HAN--PDQN uses less terminal energy per successful
task across all evaluated populations.
Its value increases from 0.948~J at 20 users to
1.491~J at 40 users, while PDQN increases from
1.040 to 1.537~J.
At 30 users, HAN--PDQN requires 1.242~J per successful
task, compared with 1.318~J for PDQN and 2.376~J
for Min-Distance.
The combination of lower delay, fewer deadline violations,
and lower energy per successful task indicates an
improvement across the three task-service dimensions.

Full-Local exhibits nearly constant delay and energy
per successful task as the population changes, consistent
with users operating independent local processors.
Its deadline violation rate nevertheless remains
approximately 49\%, and its energy per successful
task is approximately 4.1~J.
The joint consideration of delay and task outcomes
therefore reveals the service limitations of relying
entirely on local computation.

For an aggregate comparison, we first take the arithmetic
mean of each metric over the five user populations
and then compute its relative reduction against PDQN.
Under this convention, HAN--PDQN reduces successful-task
delay, deadline violation rate, and energy per successful
task by 2.16\%, 21.26\%, and 5.72\%, respectively.
The largest relative improvement occurs in deadline
compliance, while the delay and energy improvements
remain consistent across the evaluated populations.



\subsection{Service Continuity and Operating Costs}
\label{subsec:service_operating_costs}

Fig.~\ref{fig:service_operating_costs} complements the
task-level results by examining service interruption,
aggregate terminal energy, and handover activity.
These metrics reveal how different association and
offloading policies balance connection availability
against the conditions under which tasks are executed.

In Fig.~\ref{fig:service_operating_costs}(a), HAN--PDQN
has lower mean blocked user-time than PDQN at every
population. At 30 users, the values are 284.08 and
298.12~user-s per episode, respectively.
The cumulative quantity increases with population,
partly because more users contribute to the total
observation time.
Min-Distance yields a slightly lower value of
276.32~user-s at 30 users, but its deadline violation
rate in Fig.~\ref{fig:task_performance}(b) is much higher.
Connection availability therefore needs to be considered
together with the processing conditions and workload
assigned to the selected satellite.

Fig.~\ref{fig:service_operating_costs}(b) shows that
HAN--PDQN also maintains the lowest mean total
terminal energy over the evaluated populations.
At 30 users, it consumes 6.18~kJ per episode,
compared with 6.37~kJ for PDQN,
8.41~kJ for Min-Distance, and 11.26~kJ for Full-Local.
Total energy measures the system's accumulated
terminal expenditure, whereas
Fig.~\ref{fig:task_performance}(c) measures the
expenditure per successful task.
The improvement in both quantities shows that the
energy-efficiency gain accompanies a reduction in
aggregate terminal consumption.

The handover results in
Fig.~\ref{fig:service_operating_costs}(c) show a
different ordering among the methods.
At 30 users, HAN--PDQN performs 0.606 handovers per
user-minute, compared with 0.659 for PDQN,
0.192 for HAN+MAPPO, and 0.198 for Min-Distance.
Thus, HAN--PDQN obtains its task-performance advantage
with fewer handovers than PDQN, while remaining more
active in satellite reassociation than HAN+MAPPO
and Min-Distance.
The lower switching activity of the latter policies
is accompanied by higher task delay, deadline violations,
and energy per successful task.
This comparison reflects the role of the long-term
objective in balancing immediate switching costs
against subsequent task-service conditions.

Using the same across-population averaging convention
as in Section~\ref{subsec:user_scaling_performance},
HAN--PDQN reduces mean blocked user-time,
total terminal energy, and handover frequency relative
to PDQN by 7.85\%, 3.19\%, and 14.82\%, respectively.
These operating-cost reductions accompany the
task-level improvements, linking the quality of
completed service to the association and workload
decisions that sustain it.

\FloatBarrier

\section{Conclusion}
\label{sec:conclusion}

In this paper, we investigated joint handover and partial
task offloading in multi-user LEO satellite edge networks.
We formulated satellite selection and continuous workload
allocation as a sequential decision problem and developed
HAN--PDQN to account for dynamic connectivity and shared
resources. HAN distinguishes established contention through
serving relations from potential contention through common
satellite alternatives. The resulting representation conditions
both the offloading fraction generated for each candidate
and the long-term value of the corresponding satellite--fraction
pair. This design connects the destination and amount of
offloaded work through the same multi-user contention context.

Simulations over five training seeds and user populations
of 20--40 showed that, compared with PDQN, HAN--PDQN
reduced average successful-task delay, deadline violation
rate, and energy per successful task by 2.16\%, 21.26\%,
and 5.72\%, respectively. These task-level improvements
were accompanied by reductions in blocked user-time,
total terminal energy, and handover frequency. The results
demonstrate the effectiveness of the proposed framework
in balancing task reliability, processing delay, and terminal
energy across the evaluated user populations. Future work
will extend the framework to jointly allocate communication
and computing resources and investigate policy transfer
across constellation and traffic configurations.

\bibliographystyle{IEEEtran}
\bibliography{sample}
\end{document}

